%------------------------------------------------------------------------------%
\begin{question}
  Calcule o produto dos números complexos

  \[
    z_1 = 2\left(\cos\left(\frac{\pi}{6}\right) + i\sin\left(\frac{\pi}{6}\right)\right)
  \]

  \[
    z_2 = 5\left(\cos\left(\frac{\pi}{3}\right) + i\sin\left(\frac{\pi}{3}\right)\right)
  \]
\end{question}

%------------------------------------------------------------------------------%
\begin{resolution}{Solução}
  \(z_1 = 2\left(\cos\left(\frac{\pi}{6}\right) + i\sin\left(\frac{\pi}{6}\right)\right)\) e
  \(z_2 = 5\left(\cos\left(\frac{\pi}{3}\right) + i\sin\left(\frac{\pi}{3}\right)\right)\)

  \[
    \pause \rho_1
    \pause = \abs{z_1}
    \pause = 2
    \qquad\qquad
    \pause \varphi_1
    \pause = \arg(z_1)
    \pause = \frac{\pi}{6}
  \]
  \[
    \pause \rho_2
    \pause = \abs{z_2}
    \pause = 5
    \qquad\qquad
    \pause \varphi_2
    \pause = \arg(z_2)
    \pause = \frac{\pi}{3}
  \]

  \[
    \pause \rho
    \pause = \abs{z_1z_2}
    \pause = \rho_1 \rho_2
    \pause = 2 \times 5
    \pause = 10
  \]
  \[
    \pause \varphi
    \pause = \arg(z_1z_2)
    \pause = \varphi_1 + \varphi_2
    \pause = \frac{\pi}{6} + \frac{\pi}{3}
    \pause = \frac{\pi + 2\pi}{6}
    \pause = \frac{3\pi}{6}
    \pause = \frac{\pi}{2}
  \]

  \[
    \pause z
    \pause = z_1 z_2
    \pause = \rho\left(\cos(\varphi)+i\sin(\varphi)\right)
    \pause = 10\left(\cos\left(\frac{\pi}{2}\right)+i\sin\left(\frac{\pi}{2}\right)\right)
    \pause = 10i
  \]
\end{resolution}

%------------------------------------------------------------------------------%
